% IEEE conference paper: AGE — an autonomous engineering agent
% Compile: pdflatex ieee_manuscript.tex  (figures in ../paper/fig)
\documentclass[conference,10pt]{IEEEtran}

\usepackage{graphicx}
\usepackage{amsmath,amssymb}
\usepackage{booktabs}
\usepackage[table]{xcolor}
\usepackage{url}
\usepackage[margin=1in]{geometry}
\ifCLASSINFOpdf
\else
\fi

\hyphenation{Phys-For-mer}

\begin{document}

\title{AGE: An Autonomous Engineering Agent with Physics-Informed Transformers and Verified Trial-and-Error Learning}

\author{\IEEEauthorblockN{Sehaj Randhir Singh}
\IEEEauthorblockA{\textit{Department of Electrical and Computer Engineering,}\\
\textit{NYU Tandon School of Engineering (independent researcher)}\\
Brooklyn, NY, USA\\
sehajrsinghs@gmail.com}}

\maketitle

\begin{abstract}
Engineering is an interactive discipline: a design is only correct once it has
been verified against tests, simulations, or measurements. We present AGE, a
complete autonomous engineering agent that learns through trial and error
gated by verification. A planner (deterministic ``mechanical'' brain or an
LLM brain) selects actions from a sandboxed skill layer; every critical step
is verified against tests or numerically \emph{independent} simulations
before it is accepted; failures are journaled as lessons and replayed into
future plans. For physics, AGE introduces PhysFormer, a transformer whose
reasoning layers are trained jointly with physics-consistency layers that
penalize violations of the governing equations (a PINN-style residual).
Across eight physics domains (projectile, pendulum, spring--mass, simply
supported beam, cantilever, viscous Burgers, RC circuit, and 2D steady-state
heat conduction), closed-form answers are verified to residuals below
$5\times10^{-2}$ (often $10^{-9}$) by independent integrators. A 75-run
baseline study (5 domains $\times$ 3 model kinds $\times$ 5 seeds) with a
pooled Wilcoxon signed-rank test shows the physics term significantly
reduces the predicted trajectories' violation of the governing equations
($p<0.001$, Cliff's $\delta=-0.92$, $\sim$19$\times$ median reduction) while
leaving scalar accuracy statistically unchanged; an external per-instance
DeepXDE PINN baseline confirms the comparison is not self-referential. A
verification-gate benchmark (7 missions, gate on/off) shows the gate
eliminates false successes (0\% vs.\ 29\%). The complete system passes 42
unit tests and completes end-to-end design missions in a single planning
iteration, in roughly 5,900 lines of dependency-light code.
\end{abstract}

\begin{IEEEkeywords}
autonomous agents; physics-informed machine learning; transformers;
verification; engineering automation
\end{IEEEkeywords}

\section{Introduction}
\IEEEPARstart{M}{odern} machine learning excels at prediction: given a static
input, produce a static output---an image label, a sentence, a scalar
property~\cite{lecun2015,brown2020}. Engineering, by contrast, is an
\emph{interactive} discipline. An engineer does not merely predict whether a
beam will fail; she designs it, simulates it, builds it, measures it, and
revises it when the measurement disagrees with the model. The defining loop
of engineering is therefore closed-loop trial and error in which the
\emph{world}---a test suite, a simulator, or a physical instrument---is the
final arbiter of correctness~\cite{sutton2019}.

Commercial efforts such as Prometheus (founded by Jeff Bezos and Vik Bajaj,
\$12B raised at a \$41B valuation, June 2026) are pursuing this vision as an
``artificial general engineer'' for physical products~\cite{techcrunch2026}.
Yet no complete, minimal, and verifiable instantiation of the concept has
been made public: the agent literature verifies weakly, and the
physics-informed machine learning literature predicts but does not
act~\cite{raissi2019,karniadakis2021,lu2021}.

This paper makes the concept concrete. We present \textbf{AGE}, an open
autonomous engineering agent built from four components: (i)~a planner
(deterministic or language-model driven); (ii)~a sandboxed skill layer
(inspect, read, search, edit, run, verify); (iii)~a \emph{verification gate}
that accepts a step only if it passes tests or an independent numeric
simulation; and (iv)~an episodic journal that converts failed verifications
into lessons for future plans. The physics core introduces
\textbf{PhysFormer}, a transformer whose reasoning layers are trained jointly
with physics-consistency layers---a physics-informed, PINN-style
architecture in which the predicted trajectory itself is penalized for
violating the governing equations.

Our contributions are: (1)~a minimal, auditable architecture for verified
trial-and-error engineering, demonstrated end-to-end on software and physics
tasks; (2)~PhysFormer, a physics-adjusted transformer with a scale-free
physics gradient and a shape/scale trajectory decomposition, evaluated on
eight physics domains; (3)~a rigorous empirical study---75-run multi-seed
baseline matrix, pooled significance testing, and an external DeepXDE PINN
comparison---showing that the physics term's robust benefit is
\emph{trajectory-level physical consistency}, not scalar accuracy; and
(4)~a controlled benchmark showing that a verification gate eliminates false
successes in an autonomous loop.

\section{Related Work}
\textbf{LLM agents with tools and verification.} ReAct~\cite{yao2023} and
Reflexion~\cite{shinn2023} established interleaved reasoning/acting and
verbal self-reflection; tool-use and web agents~\cite{wang2023,park2023}
extend them. Recent coding agents~\cite{jimenez2024} ground success in test
execution, and verifier-based loops~\cite{li2022} use learned models to
judge outputs. AGE differs in making verification a \emph{hard gate} rather
than a soft signal, and in anchoring it to numerically independent physics
simulators, not just test suites.

\textbf{Physics-informed machine learning.} PINNs~\cite{raissi2019} embed
PDE residuals into the loss; physics-informed learning more broadly is
surveyed in~\cite{karniadakis2021}; DeepXDE~\cite{lu2021} is the standard
implementation used as an external baseline here. Gradient pathologies in
PINN training~\cite{wang2021} motivate our scale-free residual normalization.
PDEBench~\cite{takamoto2022} provides standard benchmarks. AGE's contribution
is not a new PINN but the integration of a physics-consistency layer into an
\emph{acting, verifying agent}.

\textbf{Autonomous engineering.} Design automation has a long history in
mechanics and electronics, but closed-loop agents that design, verify, and
learn from failures are new. To our knowledge AGE is the first complete,
public, minimal system that combines an agent loop, physics-informed
learning, and independent numeric verification.

\section{System Design}
\subsection{The agent loop}
The agent executes a plan of skill invocations inside a sandboxed working
directory. Each step is annotated with an expected outcome
(\texttt{expect: ok}); a gated step that fails ends the mission, writes a
lesson to the episodic journal, and the lesson is replayed into future
plans. The planner is either a deterministic ``mechanical'' brain (a
task-scripted executor, zero dependencies, fully unit-testable) or an LLM
brain (free-form planning through the same skills). Figure~\ref{fig:arch}
shows the architecture.

\begin{figure}[t]
\centering
\includegraphics[width=\linewidth]{../paper/fig/fig1_architecture.png}
\caption{AGE system architecture: planner $\rightarrow$ skills $\rightarrow$
sandboxed environment, with a verification gate and an episodic journal that
feeds lessons back into planning.}
\label{fig:arch}
\end{figure}

\subsection{Sandboxing and safety}
Every file operation is path-confined to the working directory. The shell
skill enforces a command denylist (\texttt{rm -rf /}, \texttt{git push},
\texttt{git reset --hard}, privilege escalation, disk formatting) unless an
explicit environment variable is set. Python bytecode caches are purged
before every \texttt{python} command so the verification gate always
recompiles from source---a same-length, same-second edit can otherwise be
masked by a stale \texttt{.pyc}, a failure mode we observed and fixed.

\subsection{The verification gate}
Every critical step declares an expected outcome. For software tasks the
gate runs the project's test suite; for physics tasks it invokes
\texttt{physx}, whose answers are cross-checked by numerically independent
simulators (Section~\ref{sec:physics}). A gated failure is unambiguous: the
mission fails, and the lesson is journaled. This converts ``the model
believes it succeeded'' into ``the world confirms it succeeded.''

\section{The Physics Core}
\label{sec:physics}
\subsection{Domains and independent verification}
The physics layer (\texttt{physx}) implements eight domains---projectile
motion, pendulum, spring--mass oscillator, simply supported beam, cantilever
beam, viscous Burgers conservation law, RC circuit, and 2D steady-state heat
conduction---each with three artifacts: an exact closed-form solution (the
training target), a reference trajectory sampled from it (the data), and an
\emph{independent} numeric simulator used for verification. Crucially, the
verifiers never reference the closed forms: Euler/RK4 integration for the
dynamical domains, finite-difference solution of the fourth-order beam
equation, a finite-volume upwind solver for Burgers, and a sparse
finite-difference Poisson solve for the heat plate. Verification therefore
checks the agent against physics, not against its own formulas. Closed-form
answers agree with the independent numerics to residuals below
$5\times10^{-2}$ in all domains (often $10^{-9}$ for the bending problems).

\subsection{PhysFormer: reasoning layers and physics layers together}
PhysFormer (Figure~\ref{fig:physformer}) maps a problem's parameters to both
a scalar answer and a trajectory. Parameters are $z$-scored and presented as
token pairs (parameter-id, value); a transformer encoder (48-dim, 4 heads, 3
layers) processes them; two heads produce the answer and the trajectory
(50--400 points of the dependent variable on the \emph{true} grid). The
physics-consistency layer computes how far the predicted trajectory strays
from the governing equation and adds this residual to the loss with a
warm-up schedule:
\begin{equation}
\mathcal{L} = \underbrace{\|a-a^{*}\|^{2}}_{\text{answer}}
+ w_{t}\underbrace{\|T-T^{*}\|^{2}}_{\text{trajectory data}}
+ w_{p}(t)\underbrace{\frac{\rho}{\rho_{\text{det}}+\epsilon}}_{\text{scale-free physics}},
\label{eq:loss}
\end{equation}
where $\rho$ is the mean governing-equation residual of the de-scaled
predicted trajectory and $\rho_{\text{det}}$ its detached magnitude, making
the physics gradient scale-free (a remedy for the gradient pathologies
studied in~\cite{wang2021}). For domains whose trajectories span orders of
magnitude, each trajectory is normalized by its own peak (shape-norm) and
the head output passes through a sigmoid, constraining shapes to $(0,1)$ and
removing unphysical negative predictions.

\subsubsection{Law-Conditioned Attention (LCA): the equation as input}
A loss term is a weak channel for physics: it cannot tell the network
\emph{which} law governs a problem before the output exists. LCA makes the
governing equation itself part of the input. Each law is tokenized into a
fixed 22-operator vocabulary (\texttt{first\_time\_deriv},
\texttt{second\_space\_deriv}, \texttt{self\_advection},
\texttt{moment\_curvature}, \texttt{energy\_conservation},
\texttt{exponential\_relaxation}, \ldots), embedded by an MLP into a law
vector, and injected into every transformer layer as a cross-attention
key/value stream plus a scale/bias gate on the shared readout. Input tokens
are labeled by physical quantity from a shared 13-token vocabulary, so the
simply supported beam and the cantilever present \emph{identical} parameter
tokens (length, force, modulus, inertia, thickness) and only the law
signature disambiguates them. One shared transformer (shared body and shared
heads) then serves six laws at once. In a six-seed paired test against a
constant-signature ablation (identical architecture, data, and budget), the
equation signal significantly improves trajectory fidelity ($p=0.0003$ over
36 runs, 21\% median curve-error reduction) and cuts held-out answer error
on the laws whose parameters do not identify them (beam 34\% vs 40\%,
cantilever 28\% vs 35\%); an honest negative on RC, whose normalized
trajectory is identical for every time constant, localizes where
conditioning does not pay. The generalist matches single-law specialists
within 1.1--2.8$\times$ on five of six laws at 37\% of their data per law.

The equation channel is causal, not decorative. On held-out beam and
cantilever problems whose parameter tokens are literally identical, swapping
the law signature at inference displaces the model's prediction across the
boundary between the two bending solutions (steering index $+0.090$ vs.
$-0.132$ for a constant-signature control, $p<0.0001$; the control is
provably insensitive), and the generalist adapts to a law held out of
pretraining with 25\% of the specialist's data ($2.4$--$2.9\times$ lower
answer error than a from-scratch specialist on the same samples, $0.20$
vs.\ $0.58$ median). A controlled physics-vs-data-only sweep on the 2D
heat plate (identical networks, budgets, seeds; only the residual weight
differs) shows the residual-loss channel cuts the predicted field's
governing-equation violation 4--8$\times$ without improving held-out field
accuracy---consistency, not accuracy, is what the loss channel buys, and
field fidelity is carried by the data and the LCA input channel.

\begin{figure}[t]
\centering
\includegraphics[width=\linewidth]{../paper/fig/fig2_physformer.png}
\caption{PhysFormer: parameter tokens $\rightarrow$ transformer encoder
$\rightarrow$ answer head + trajectory head; the physics-consistency layer
evaluates the governing-equation residual of the de-scaled trajectory on the
true grid and feeds it back into the loss.}
\label{fig:physformer}
\end{figure}

\section{Experiments}
\subsection{Protocol}
Each domain samples uniform parameter ranges (e.g., beam:
$L\in[1,6]$\,m, $P\in[10^{3},5\times10^{4}]$\,N; Burgers:
$\nu\in[0.02,0.1]$, $A\in[0.6,2]$, $\sigma\in[0.2,0.5]$). Models train on
256 samples at 60 epochs (Burgers/heat2d headlines: 120--240 epochs with
trajectory weight 3); the seed controls data generation, shuffling, and
weight initialization, so multi-seed runs are proper statistical
replicates. Validation is the last $\sim$1/8 of the generated problems;
held-out figures use 250 freshly generated problems per domain.

\subsection{Baselines: attention and physics, isolated and measured}
We trained three model kinds on identical budgets---an MLP encoder, the
transformer without the physics term, and PhysFormer with it---across five
domains $\times$ five seeds (75 runs). Table~\ref{tab:baselines} reports
mean$\pm$std; Fig.~\ref{fig:baselines} visualizes them. Two findings stand
out.

\emph{The physics term buys trajectory consistency.} Physics-on models
violate the governing equations by roughly an order of magnitude less than
the no-physics transformer (beam residual $3.4$ vs.\ $48$; cantilever $39$
vs.\ $556$; Burgers $0.01$ vs.\ $0.09$; heat2d $0.45$ vs.\ $10.4$).

\emph{Are these differences significant?} With five seeds per kind, the
per-domain Wilcoxon signed-rank test is limited to $p\geq0.0625$, so we pool
all five domains into 25 paired runs per comparison. On the trajectory
residual, physics-on vs.\ no-physics gives $p<0.001$ with Cliff's
$\delta=-0.92$ (large effect) and a $\sim$19-fold median reduction;
physics-on vs.\ MLP gives $p<0.001$, $\delta=-0.84$. On the scalar answer,
no comparison reaches significance ($p\geq0.65$, $|\delta|\leq0.12$). The
physics term therefore does \emph{not} move the scalar answer---it makes the
full trajectory physically consistent, significantly and with large effect.

\emph{Scalar accuracy is domain-dependent.} PhysFormer matches or beats the
MLP where the parameter space is rich (beam $14.7\%$ vs.\ $15.6\%$;
cantilever $14.8\%$ vs.\ $15.2\%$; Burgers $0.9\%$ vs.\ $1.8\%$; heat2d
$1.5\%$ vs.\ $1.6\%$ held-out) but loses on the trivial two-parameter
projectile ($10.8\%$ vs.\ $6.5\%$): an honest negative---attention and
physics pay for themselves only when the map and trajectory are rich enough.

\begin{table}[t]
\caption{Equal-budget baselines (256 samples, 60 epochs, five seeds,
mean$\pm$std). Left: validation relative error of the scalar answer. Right:
trajectory residual (mean governing-equation violation of the predicted
trajectory).}
\label{tab:baselines}
\centering
\small
\begin{tabular}{lcccc}
\toprule
 & \multicolumn{2}{c}{Answer error (\%)} & \multicolumn{2}{c}{Trajectory residual} \\
\cmidrule(lr){2-3}\cmidrule(lr){4-5}
Domain & MLP & No-phys & $-$physics & $+$physics \\
\midrule
Beam & $15.6\pm2.3$ & $19.3\pm7.0$ & $48\pm14$ & $3.4\pm0.5$ \\
Cantilever & $15.2\pm1.5$ & $20.4\pm7.5$ & $556\pm214$ & $39\pm12$ \\
Projectile & $6.5\pm1.9$ & $7.6\pm5.6$ & $0.12$ & $0.10$ \\
Burgers & $1.8\pm0.7$ & $1.0\pm0.3$ & $0.09$ & $0.01$ \\
Heat2D & $1.6\pm0.4$ & $1.4\pm0.5$ & $10.4\pm5.9$ & $0.45\pm0.02$ \\
\bottomrule
\end{tabular}
\end{table}

\begin{figure}[t]
\centering
\includegraphics[width=\linewidth]{../paper/fig/fig9_baselines.png}
\caption{Multi-seed baselines (mean$\pm$std over five seeds): MLP, no-physics
transformer, and transformer $+$ physics. The physics term's robust benefit
is trajectory consistency.}
\label{fig:baselines}
\end{figure}

\subsection{Held-out generalization and trajectory fidelity}
On 250 held-out problems per domain, PhysFormer achieves mean relative
errors of $14.8\%$ (beam), $11.6\%$ (cantilever), $9.6\%$ (projectile),
$\sim$1\% (Burgers peak velocity), and $1.9\%$ (heat2d peak temperature),
with $R^{2}\geq0.99$ on the wide-range domains (Fig.~\ref{fig:heldout}).
The shape/scale decomposition reproduces full trajectories: the beam
deflection curve within $\sim$19\% of peak, the cantilever within
$\sim$16\%, the Burgers final-time field within $\sim$11\%, and the heat2d
field within $\sim$6\% of peak (mean error $\sim$1.3\%) at a 240-epoch
trajectory budget---the 400-point oscillatory field needs roughly twice
the epochs of the one-dimensional domains (the 60-epoch model underfits it
with $\sim$50\% peak error). The sigmoid constraint removes unphysical
negative predictions.

\begin{figure}[t]
\centering
\includegraphics[width=\linewidth]{../paper/fig/fig4_predictions.png}
\caption{Held-out predictions (250 problems per domain) for beam,
cantilever, projectile, and heat2d.}
\label{fig:heldout}
\end{figure}

\subsection{Nonlinear and 2D domains}
The viscous Burgers equation $u_{t}+uu_{x}=\nu u_{xx}$ (nonlinear
conservation law with shock formation) stresses every design choice: the
predicted object is a full $25\times8$ spatiotemporal field, the exact
reference comes from the Cole--Hopf transform, and the independent verifier
is a finite-volume upwind solver. The physics-on models reach a trajectory
residual of $0.01$, an order of magnitude below no-physics ($0.09$), and
held-out peak-velocity error is $\sim$1\%; the canonical shock profile is
resolved to $\sim$11\% of peak. The 2D heat plate (Poisson on a
$20\times20$ grid) shows the largest relative residual improvement of any
domain ($0.45$ vs.\ $10.4$, $23\times$): the Laplacian constraint couples
every interior point to its neighbors and cannot be satisfied by chance on a
400-point field (Fig.~\ref{fig:heat2d}).

\begin{figure}[t]
\centering
\includegraphics[width=\linewidth]{../paper/fig/fig10_heat2d.png}
\caption{Steady-state heat conduction: exact vs.\ predicted temperature
field, pointwise error, and held-out peak temperatures.}
\label{fig:heat2d}
\end{figure}

\subsection{External PINN baseline (DeepXDE)}
To ensure the comparison is not self-referential, we trained the standard
PINN library DeepXDE~\cite{lu2021} (v1.15.0, PyTorch backend) on three
representative Burgers problems, giving it the setup most favourable to a
PINN: \emph{one fresh network per problem instance} (2--50$\times$4--1 tanh
FNN, 2{,}000 collocation points, Adam 12k + L-BFGS 2k steps, exact
Cole--Hopf boundary data). PhysFormer is a single network trained once on
256 instances that generalizes to unseen parameters. Both are evaluated with
identical metrics on the model grid (Fig.~\ref{fig:deepxde}).

On all three instances the per-instance PINN converges to essentially exact
solutions (final-time peak error $\leq0.01\%$; full-field error
0.03--0.06\%; $\sim$19--21 minutes of training per instance), which
validates both the problem setup and the Cole--Hopf ground truth.
PhysFormer---one network, one $\sim$minute training, then a millisecond
forward pass for any instance---achieves final-time peak errors of
0.02--3.7\% and full-field errors of 33--51\% on the same instances:
worse on any given problem, as expected, but with no per-instance
optimization at all. Notably, PhysFormer's governing-equation residual on
the predicted field is comparable to or lower than the converged PINN's
($2\times10^{-4}$ vs.\ $1$--$5\times10^{-3}$): the generalist's
predictions are as physical as the specialist's, just less accurate. The
comparison quantifies the generalist--specialist tradeoff directly.

\begin{figure}[t]
\centering
\includegraphics[width=\linewidth]{../paper/fig/fig11_deepxde.png}
\caption{Per-instance DeepXDE PINN vs.\ PhysFormer on three Burgers
instances (canonical, mild, shock): answer error, full-field error, and
governing-equation residual.}
\label{fig:deepxde}
\end{figure}

\subsection{The verification gate}
We ran seven deterministic missions---benign software tasks, beam,
cantilever, and Burgers design tasks, an injected-fault task, a skill-level
guard, and a TODO scan---twice each: with the gate enforced and with it
disabled (the verifier-free baseline). Ground truth was computed
independently of the agent. Table~\ref{tab:gate} shows the result: the gate
produces \emph{zero} false successes, while the verifier-free agent claims
success on 29\% of missions that actually fail. The gate converts claimed
success into verified success.

\begin{table}[t]
\caption{Verification-gate benchmark (7 missions). False success = the agent
reports success but independent ground truth fails.}
\label{tab:gate}
\centering
\small
\begin{tabular}{lcc}
\toprule
 & Gate enforced & Gate disabled \\
\midrule
True successes & 5/7 & 5/7 \\
False successes & \textbf{0/7} & 2/7 \\
False-success rate & \textbf{0\%} & \textbf{29\%} \\
\bottomrule
\end{tabular}
\end{table}

\section{Discussion and Limitations}
The paper's central empirical claim is now precisely stated: physics
supervision inside the transformer does not improve scalar accuracy
(no significant difference, $p=0.65$ pooled) but does improve
\emph{trajectory physicality} significantly ($p<0.001$, $\delta=-0.92$),
and this holds across linear, nonlinear, and 2D domains and against an
external PINN implementation.

Limitations are explicit: (1)~\emph{Scale}: 3--7 parameters per domain,
CPU-minute training, 1D and one 2D field; generalization errors of
$\sim$10--15\% at the 256-sample budget. (2)~\emph{Trajectory fidelity}:
smooth profiles are reproduced well but sharp features remain coarse---the
Burgers shock case shows $\sim$52\% curve error on the $25\times8$ grid
even though the scalar is within 5\%. (3)~\emph{Repair}: the mechanical
brain inspects and solves but does not yet author fixes for failing tests.
(4)~\emph{Simulated physics}: verification is against numeric integration,
not sensors or manufactured parts; the architecture ports unchanged by
swapping the verify skill for a measurement harness. (5)~\emph{Physics
classes}: quasi-static bending and scalar conservation laws, not general
continuum mechanics or multi-physics coupling.

\section{Conclusion}
We presented AGE, a complete autonomous engineering agent that learns
through verified trial and error, and PhysFormer, a transformer adjusted
for physics whose reasoning and physics-consistency layers are trained
together. A 75-run multi-seed study with pooled significance testing, an
external DeepXDE baseline, and a controlled gate benchmark support three
transferable claims: verification, not model confidence, must arbitrate
success; physics and data should be trained jointly with a scale-free
physics gradient; and memory should be episodic and replayable. The system
is roughly 5,900 lines of dependency-light code, passes 54 unit tests, and
runs on a laptop---a fully auditable baseline for the ``general engineer''
agenda that commercial efforts are pursuing at much larger scale.

\begin{thebibliography}{99}
\bibitem{lecun2015} Y.~LeCun, Y.~Bengio, and G.~Hinton, ``Deep learning,''
\emph{Nature}, vol.~521, pp.~436--444, 2015.
\bibitem{brown2020} T.~Brown \emph{et al.}, ``Language models are few-shot
learners,'' in \emph{Advances in Neural Information Processing Systems 33},
2020.
\bibitem{sutton2019} R.~Sutton, ``The bitter lesson,''
\url{http://www.incompleteideas.net/IncIdeas/BitterLesson.html}, 2019.
\bibitem{techcrunch2026} M.~Temkin, ``Jeff Bezos's Prometheus raises \$12B
to build an `artificial general engineer' for the physical world,''
\emph{TechCrunch}, Jun.~2026.
\bibitem{raissi2019} M.~Raissi, P.~Perdikaris, and G.~E. Karniadakis,
``Physics-informed neural networks: a deep learning framework for solving
forward and inverse problems involving nonlinear partial differential
equations,'' \emph{J. Comput. Phys.}, vol.~378, pp.~686--707, 2019.
\bibitem{karniadakis2021} G.~E. Karniadakis \emph{et al.}, ``Physics-informed
machine learning,'' \emph{Nat. Rev. Phys.}, vol.~3, pp.~422--440, 2021.
\bibitem{lu2021} L.~Lu, X.~Meng, Z.~Mao, and G.~E. Karniadakis,
``DeepXDE: a deep learning library for solving differential equations,''
\emph{SIAM Rev.}, vol.~63, pp.~208--228, 2021.
\bibitem{wang2021} S.~Wang, Y.~Teng, and P.~Perdikaris, ``Understanding and
mitigating gradient flow pathologies in physics-informed neural networks,''
\emph{SIAM J. Sci. Comput.}, vol.~43, pp.~A3055--A3081, 2021.
\bibitem{takamoto2022} M.~Takamoto \emph{et al.}, ``PDEBench: an extensive
benchmark for scientific machine learning,'' in \emph{Advances in Neural
Information Processing Systems 35}, 2022.
\bibitem{yao2023} S.~Yao \emph{et al.}, ``ReAct: synergizing reasoning and
acting in language models,'' in \emph{ICLR}, 2023.
\bibitem{shinn2023} N.~Shinn \emph{et al.}, ``Reflexion: language agents
with verbal reinforcement learning,'' in \emph{NeurIPS}, 2023.
\bibitem{wang2023} G.~Wang \emph{et al.}, ``Voyager: an open-ended embodied
agent with large language models,'' \emph{arXiv:2305.16291}, 2023.
\bibitem{park2023} J.~S. Park \emph{et al.}, ``Generative agents:
interactive simulacra of human behavior,'' in \emph{UIST}, 2023.
\bibitem{jimenez2024} C.~E. Jimenez \emph{et al.}, ``SWE-bench: can language
models resolve real-world GitHub issues?,'' in \emph{ICLR}, 2024.
\bibitem{li2022} Y.~Li \emph{et al.}, ``Training verifiers to solve math
word problems,'' \emph{arXiv:2110.14168}, 2022.
\end{thebibliography}

\end{document}
